% Tema 1: límites y continuidad.
%
% Un documento de Didacta es una portada y una lista de lecciones. Nada más:
% ni preámbulo de sesenta líneas, ni \import con rutas ../../../, ni una copia
% del contenido. Las lecciones están en la biblioteca y aquí solo se nombran.
%
% El mismo fichero da todas las salidas del tema. Desde la aplicación se elige
% cuál con el botón de compilar; en el terminal,
%
%   didacta build tema-1                   todas las de year.yaml
%   didacta build tema-1 -p slides -l va   las diapositivas, en valenciano

\input{didacta-bootstrap}
\usepackage{didacta}

%% El título y los datos de la asignatura salen de `course.yaml` y
%% `year.yaml`, en el idioma que se está compilando: Didacta los inyecta al
%% compilar. Lo de aquí abajo es solo un respaldo, para que este fichero
%% compilado fuera de Didacta siga teniendo una portada con título.
\DidactaCourse{
  title   = {Cálculo I},
  teacher = {Tu nombre},
}
\DidactaDocument{Tema 1: límites y continuidad}

\begin{document}
\DidactaTitlePage
\DidactaContents

%% A partir de aquí va la composición, y la escribe el motor a partir de
%% `year.yaml`, que es quien manda: lo que se añada, se quite o se reordene
%% desde la aplicación se pone al día aquí al compilar. Una línea que empieza
%% por `%%` es una lección comentada, que este año no se da.

\DidactaSection{es={Límite de una función}, va={Límit d'una funció}, en={The limit of a function}}
\DidactaUnit{u-240b481e5d63}
\DidactaUnit{u-6431387efe7b}
%% \DidactaUnit{u-e08461a36e86}
\DidactaUnit{u-2d7a8345e536}

\DidactaSection{es={Continuidad}, va={Continuïtat}, en={Continuity}}
\DidactaUnit{u-a1bdfa8905d7}
\DidactaUnit{u-245a446424c6}

\end{document}
